\clearpage
\section{Schur--Weyl quotients and the boundary of fusion}
\label{sec:f1-limit-schur}

At the symmetric endpoint, the Hecke tower becomes the symmetric-group
tower.  Its action on tensor powers of $\mathbb C^d$ gives a concrete
monoidal quotient, but it changes the physical trace: a permutation is
weighted by its number of cycles rather than only by whether it is the
identity.  The difference disappears at fixed tensor degree as $d$ grows,
and the associated conditional expectations converge with an explicit
finite-$d$ bound.

Rigidity and fusion require further data.  The nonnegative grading of the
universal amplitude category rules out evaluation and coevaluation maps
inside that category.  A rigid comparison must name a target, a monoidal
functor, dual objects, cups and caps, dagger and trace compatibility.  A
Temperley--Lieb or Fibonacci branch adds quotient ideals and a different
trace; neither is an automatic consequence of the positive-real endpoint.

\subsection{The tensor-power quotient}

\begin{definition}[The symmetric endpoint amplitude category]
For real $q>0$, let $H_n(q)$ be the positive type-$A$ Hecke algebra with
self-adjoint generators $T_i$, relations
\[
 (T_i-q)(T_i+1)=0,\qquad
 T_iT_{i+1}T_i=T_{i+1}T_iT_{i+1},\qquad
 T_iT_j=T_jT_i\quad(|i-j|>1),
\]
and coefficient trace $\tau_n(T_w)=\delta_{w,1}$.  Ordered block inclusion
maps $T_u\otimes T_v$ to $T_{u\times v}$.  The finite positive composition
category $\mathcal U_q$ consists of finite sums and self-adjoint projection
retracts of the graded regular objects; unequal degrees have zero Homs.  At
$q=1$, $H_n(1)=\mathbb C[S_n]$ and the generator $x=[1]$ satisfies
$\operatorname{End}(x^{\otimes n})=\mathbb C[S_n]$.
\end{definition}
\begin{scope}
This is an amplitude category with finite $C^*$-endomorphism algebras.  Its
nonnegative grading is retained, and no internal dual for $x$ is assumed.
\end{scope}
\provenance{D1101,D1102,D1224}{theory/sidequests/f1-limit/universal.md}{theory/checks/f1\_operational\_check.py (H13, finite corner samples)}{theory/verdicts/f1-limit-adjudication.md}

\begin{definition}[Classical Schur--Weyl quotient]
Fix $d\geq1$, put $V_d=\mathbb C^d$, and let
\[
 P_\sigma(v_1\otimes\cdots\otimes v_n)
 =v_{\sigma^{-1}(1)}\otimes\cdots\otimes v_{\sigma^{-1}(n)}.
\]
Define
\[
 \rho_{d,n}:\mathbb C[S_n]\longrightarrow
 \operatorname{End}_{U(d)}(V_d^{\otimes n}),
 \qquad \rho_{d,n}(\sigma)=P_\sigma.
\]
Its kernel is the sum of Wedderburn blocks indexed by partitions of $n$ with
more than $d$ rows.  Equivalently, for $n\geq d+1$, it is the ideal generated
by the antisymmetrizer on the first $d+1$ positions.  The map is faithful
exactly when $d\geq n$.
\end{definition}
\begin{scope}
The target is the full $U(d)$-commutant on one tensor power.  For $d<n$ the
map is a quotient, and no faithful embedding is claimed.
\end{scope}
\provenance{D1232}{theory/sidequests/f1-limit/schur-weyl.md}{theory/checks/f1\_limit\_check.py (L8, n and d at most 4)}{theory/verdicts/f1-limit-adjudication.md}

\begin{definition}[Quotient assembly]
Under the standard associator
$V_d^{\otimes(m+n)}\cong V_d^{\otimes m}\otimes V_d^{\otimes n}$, define
\[
 \nu_{m,n}\bigl(\rho_{d,m}(a)\otimes\rho_{d,n}(b)\bigr)
 =\rho_{d,m+n}\bigl(\iota_{m,n}(a\otimes b)\bigr).
\]
These are associative unital algebra maps.  On right modules, the associated
Schur--Weyl functor is
\[
 M\longmapsto M\otimes_{\mathbb C[S_n]}V_d^{\otimes n},
\]
where the tensor power has its left permutation action.
\end{definition}
\begin{scope}
Well-definedness uses compatibility of the kernels with block inclusion.
The source tensor on modules is balanced induction, not an objectwise
$S_{m+n}$-action on $M\otimes N$.
\end{scope}
\provenance{D1233}{theory/sidequests/f1-limit/schur-weyl.md}{theory/checks/f1\_limit\_check.py (L8, finite rank/trace samples; assembly by proof)}{theory/verdicts/f1-limit-adjudication.md}

\begin{theorem}[Coherent finite-rank Schur--Weyl quotient]
{\normalfont\statusProved}\quad For every $d\geq1$ and $n\geq0$,
$\rho_{d,n}$ is a surjective star-homomorphism onto the tensor-power
$U(d)$-commutant.  Its kernel is exactly the sum of partition blocks with
more than $d$ rows, equivalently the antisymmetrizer-generated ideal, and it
is faithful exactly for $d\geq n$.  The maps $\nu_{m,n}$ form a
multiplicative sequence and induce a strong monoidal functor from finite
right modules with induction tensor.
\end{theorem}
\begin{scope}
The theorem is at the symmetric endpoint.  It permits a quotient kernel and
does not identify induction with the bare vector-space tensor product of
modules over different symmetric groups.
\end{scope}
\provenance{F1-LIM-SW,D1232,D1233}{theory/sidequests/f1-limit/schur-weyl.md}{theory/checks/f1\_limit\_check.py (L8, n,d at most 4; general kernel by proof)}{theory/verdicts/f1-limit-adjudication.md}
\begin{proof}
The permutation convention gives $P_\sigma P_\pi=P_{\sigma\pi}$ and
$P_\sigma^*=P_{\sigma^{-1}}$.  Position permutations commute with the
diagonal $U(d)$ action.  Conversely, the complex span of the Lie algebra of
$U(d)$ is $\mathfrak{gl}_d(\mathbb C)$, so the unitary and classical
general-linear commutants agree; classical Schur--Weyl duality makes the map
surjective.

If $d\geq n$, apply a putative linear relation among the $P_\sigma$ to
$e_1\otimes\cdots\otimes e_n$; the resulting $n!$ words are distinct.  If
$d<n$, the $(d+1)$-strand antisymmetrizer acts through
$\bigwedge^{d+1}V_d=0$.  In the decomposition
\[
 V_d^{\otimes n}\cong
 \bigoplus_{\substack{\lambda\vdash n\\\operatorname{rows}(\lambda)\leq d}}
 L_d(\lambda)\otimes S^\lambda,
\]
the block $\operatorname{End}(S^\lambda)$ survives exactly for partitions
with at most $d$ rows.  The branching rule says precisely those excluded
blocks contain the sign representation of the embedded $S_{d+1}$, proving
the ideal description.

On elementary tensors, block permutation gives
\[
 \rho_{d,m+n}(\iota(\sigma\otimes\pi))
 =\rho_{d,m}(\sigma)\otimes\rho_{d,n}(\pi).
\]
This proves kernel compatibility and associativity.  Cancelling consecutive
extension and restriction of scalars identifies the image of
$(M\star N)$ with
$(M\otimes_{\mathbb C[S_m]}V_d^{\otimes m})\otimes
 (N\otimes_{\mathbb C[S_n]}V_d^{\otimes n})$, giving the strong monoidal
structure.
\end{proof}

\subsection{Which trace does the quotient see?}

\begin{definition}[Tensor trace and physical subgroup expectation]
Define the normalized tensor trace
\[
 \operatorname{tr}_{d,n}(a)=d^{-n}\operatorname{Tr}(\rho_{d,n}(a)).
\]
For a permutation $\sigma$, this is
\[
 \operatorname{tr}_{d,n}(\sigma)=d^{\operatorname{cyc}(\sigma)-n}.
\]
When $d\geq n$ and $K=S_m\times S_{n-m}$, the physical trace-preserving
expectation $E_{d,K}:\mathbb C[S_n]\to\mathbb C[K]$ is characterized by
\[
 \operatorname{tr}_{d,n}(k^*E_{d,K}(x))
 =\operatorname{tr}_{d,n}(k^*x),\qquad k\in K.
\]
In the group basis its coefficients solve $G_dc=b$, where
\[
 G_d(k,h)=d^{\operatorname{cyc}(k^{-1}h)-n},\qquad
 b_k=d^{\operatorname{cyc}(k^{-1}\sigma)-n}.
\]
\end{definition}
\begin{scope}
Faithfulness on the full group algebra, and therefore this Gram inversion
there, requires $d\geq n$.  On the quotient image the tensor trace remains
faithful for all $d$.
\end{scope}
\provenance{D1234}{theory/sidequests/f1-limit/schur-weyl.md}{theory/checks/f1\_limit\_check.py (L8,L9, declared finite ranks)}{theory/verdicts/f1-limit-adjudication.md}

\begin{definition}[Fixed-degree stable trace comparison]
Let $\tau_n$ be coefficient trace and let $E_{\infty,K}$ delete group-basis
coefficients outside $K$.  For fixed $n$ and
$x=\sum_\sigma a_\sigma\sigma$,
\[
 \left|\operatorname{tr}_{d,n}(x)-\tau_n(x)\right|
 \leq \frac1d\sum_{\sigma\ne1}|a_\sigma|.
\]
Put $N=|K|$.  If both $d\geq n$ and $d>N-1$, then for
$\sigma\notin K$ every coefficient of $E_{d,K}(\sigma)$ has modulus at most
\[
 \frac{1}{d-(N-1)},
\]
whereas $E_{d,K}(\sigma)=\sigma$ for $\sigma\in K$.  Thus the physical
expectation converges coefficientwise to coefficient deletion at fixed
degree.
\end{definition}
\begin{scope}
Both inequalities on $d$ are required: one for faithfulness and one for the
Neumann-series bound.  No equality between the two traces or expectations
is claimed at fixed finite $d$.
\end{scope}
\provenance{D1235}{theory/sidequests/f1-limit/schur-weyl.md}{theory/checks/f1\_limit\_check.py (L8,L9; general bound analytic)}{theory/verdicts/f1-limit-adjudication.md}

\begin{theorem}[Stable physical trace and expectation]
{\normalfont\statusProved}\quad At fixed $n$, the normalized tensor trace
converges coefficientwise to $\tau_n$.  If $K=S_m\times S_{n-m}$,
$d\geq n$, and $d>|K|-1$, then $E_{d,K}$ converges coefficientwise to
$E_{\infty,K}$ with the bound above.  For fixed $h\geq0$,
$\tau_n(h)=1$, and $0\leq e\leq1$, put
\[
 \Delta_d(y)=\frac1d\sum_{\sigma\ne1}|y_\sigma|.
\]
Whenever $\Delta_d(h)<1$, the normalized quotient density has Born
probability $\operatorname{tr}_{d,n}(he)/\operatorname{tr}_{d,n}(h)$ and
\[
 \left|\frac{\operatorname{tr}_{d,n}(he)}{\operatorname{tr}_{d,n}(h)}
       -\tau_n(he)\right|
 \leq\frac{\Delta_d(he)+\Delta_d(h)}{1-\Delta_d(h)}.
\]
\end{theorem}
\begin{scope}
All convergence is at fixed tensor degree and for fixed coefficient data.
There is no uniform claim when the degree, coefficient norms, subgroup, or
experiment varies with $d$.
\end{scope}
\provenance{F1-LIM-SWTRACE,D1234,D1235}{theory/sidequests/f1-limit/schur-weyl.md}{theory/checks/f1\_limit\_check.py (L8 for n,d at most 4; L9 finite S2-in-S3 sentinel)}{theory/verdicts/f1-limit-adjudication.md}
\begin{proof}
A tensor-basis word is fixed by $P_\sigma$ exactly when it is constant on
each cycle, so there are $d^{\operatorname{cyc}(\sigma)}$ fixed words.  A
nonidentity permutation has at most $n-1$ cycles, proving the trace bound.
Moreover
\[
 \operatorname{tr}_{d,n}(x^*x)
 =d^{-n}\operatorname{Tr}(\rho_{d,n}(x)^*\rho_{d,n}(x)),
\]
so its radical is exactly the Schur--Weyl kernel.

The conditional expectation is orthogonal projection for this trace.  Its
Gram matrix is $G_d=I+R$, with
$\lVert R\rVert_\infty\leq(N-1)/d$.  Under the two displayed hypotheses,
\[
 \lVert G_d^{-1}\rVert_\infty
 \leq\frac{1}{1-(N-1)/d}.
\]
For $\sigma\notin K$, the right-hand Gram vector has norm at most $1/d$,
which proves the coefficient bound.  The finite-rank difference is real: for
$n=3$, $K=\langle(12)\rangle$, one obtains
$E_{d,K}((23))=d^{-1}1$, while coefficient deletion gives zero.

Finally the trace estimate applied to $h$ makes its physical normalization
at least $1-\Delta_d(h)$.  Apply the same estimate to $he$ and divide by
this lower bound to obtain the stated Born inequality.
\end{proof}

\subsection{Rigid and fusion targets require additional data}

\begin{definition}[Rigid target datum]
A rigid comparison consists of an additive idempotent-complete monoidal
category $\mathcal R$, a named strong monoidal functor
$i:\mathcal U_q\to\mathcal R$, and a named left and right dual $y^\vee$ of
$y=i(x)$.  It includes
\[
 \operatorname{coev}_R:\mathbf1\to y\otimes y^\vee,
 \quad \operatorname{ev}_R:y^\vee\otimes y\to\mathbf1,
\]
\[
 \operatorname{coev}_L:\mathbf1\to y^\vee\otimes y,
 \quad \operatorname{ev}_L:y\otimes y^\vee\to\mathbf1,
\]
satisfying the four zigzags with the stated associators.  The target is
generated from $y,y^\vee$ by tensor, finite sums and retracts.  The functor
may have a kernel.  A positive dagger version additionally requires a
$C^*$-category, a star-preserving $i$, compatible dagger/pivotal data, and a
positive spherical structure.  Duals of sums and retracts are formed in the
additive Karoubi closure.
\end{definition}
\begin{scope}
Full faithfulness is additional and is never inferred from the existence of
cups and caps.  The target closure is part of the datum.
\end{scope}
\provenance{D1236}{theory/sidequests/f1-limit/fusion.md}{theory/checks/f1\_limit\_check.py (L8, finite quotient witness only)}{theory/verdicts/f1-limit-adjudication.md}

\begin{definition}[Fusion quotient datum]
A fusion quotient consists of a monoidal ideal $\mathcal I$ in a chosen
rigid extension, a compatible dagger on the quotient, a positive pivotal
trace, and the instruction to quotient and then take additive Karoubi
completion.  Any trace radical or negligible ideal is named explicitly.
The result is called fusion only after proving semisimplicity, simple unit,
finite-dimensional Hom spaces, duals, and finitely many simple isomorphism
classes.
\end{definition}
\begin{scope}
A family of degreewise algebra ideals is insufficient unless it is closed
under tensor, composition, dagger, and the chosen duality operations.
\end{scope}
\provenance{D1237}{theory/sidequests/f1-limit/fusion.md}{structural proof; no finite sample proves fusion}{theory/verdicts/f1-limit-adjudication.md}

\begin{definition}[Temperley--Lieb branch]
For $q>0$, set
\[
 f_i=\frac{q-T_i}{q+1},\qquad
 \delta=\sqrt q+\frac1{\sqrt q}.
\]
A Temperley--Lieb branch imposes the tensor-ideal relations
\[
 f_if_{i+1}f_i=\delta^{-2}f_i,\qquad
 f_{i+1}f_if_{i+1}=\delta^{-2}f_{i+1},
\]
equivalently killing the excluded three-strand Hecke summand.  It also
supplies diagrammatic cups and caps, dagger, and the Jones--Markov trace.
The faithful coefficient trace is not declared to descend through this
nonzero ideal.
\end{definition}
\begin{scope}
This is a named quotient branch with new rigid and trace data.  It is not a
property of the full positive flag/Hecke category.
\end{scope}
\provenance{D1238}{theory/sidequests/f1-limit/fusion.md}{theory/checks/f1\_operational\_check.py (H10, finite trace obstruction)}{theory/verdicts/f1-limit-adjudication.md}

\begin{definition}[Fibonacci root branch]
The unitary Fibonacci branch is the even part of the level-three
$SU(2)$/Jones fusion quotient.  Choose
\[
 q_H=e^{2\pi i/5},\qquad v=e^{\pi i/5},\qquad v^2=q_H,
 \qquad q_{\mathrm{ILZ}}=-v,
\]
and distinguish the loop generator from the projection by
\[
 U_i^{\mathrm{ILZ}}=\delta f_i,\qquad
 \delta=-(q_{\mathrm{ILZ}}+q_{\mathrm{ILZ}}^{-1})
       =v+v^{-1}=2\cos(\pi/5).
\]
Quotient by the Jones trace radical, take additive Karoubi completion, and
retain the even labels $\{0,2\}$, with
\[
 2\otimes2\cong0\oplus2.
\]
\end{definition}
\begin{scope}
The square parameter has order five.  This root-of-unity construction is
not the $q=1$ specialization of the positive-real branch, where $\delta=2$
and the spin tower is untruncated.
\end{scope}
\provenance{D1239}{theory/sidequests/f1-limit/fusion.md}{primary-source parameter derivation; no cyclotomic checker}{theory/verdicts/f1-limit-adjudication.md}

\begin{theorem}[A positive rigid target at the symmetric endpoint]
{\normalfont\statusProved}\quad The category $\mathcal U_q$ is not rigid:
its degree-one generator has no left or right dual.  A positive rigid
comparison requires the complete datum above.  At $q=1$, for every $d\geq1$,
the full replete $C^*$-tensor subcategory of $\operatorname{Rep}(U(d))$
generated by $V_d,V_d^*$, finite sums and subobjects is a concrete positive
rigid target of the Schur--Weyl quotient image.  The named functor
$i_d:\mathcal U_1\to\mathcal R_d$ may have a kernel.
\end{theorem}
\begin{scope}
The comparison is not asserted fully faithful in degrees $n>d$.  The target
has infinitely many simple objects as degrees vary and is not thereby a
fusion category.
\end{scope}
\provenance{F1-LIM-RIGID,D1236,D1237}{theory/sidequests/f1-limit/fusion.md}{theory/checks/f1\_limit\_check.py (L8, finite Schur--Weyl kernel samples)}{theory/verdicts/f1-limit-adjudication.md}
\begin{proof}
Every summand of $x\otimes y$ has degree at least one for any object $y$ of
the nonnegatively graded category, while the unit has degree zero.  Thus an
evaluation $x\otimes y\to\mathbf1$ must vanish and cannot satisfy a zigzag.

At the endpoint, send $x$ to $V_d$, symmetric-group arrows to their tensor
permutations, and a projection retract to the range of its image projection.
This is a named strong monoidal star-functor, with precisely the
Schur--Weyl kernel.  The usual pairing and copairing
\[
 \operatorname{ev}(\varphi\otimes v)=\varphi(v),\qquad
 \operatorname{coev}(1)=\sum_i e_i\otimes e_i^*
\]
satisfy the zigzags by basis contraction.  The Hilbert dagger and
$V_d\cong V_d^{\vee\vee}$ give positive dimension $d$ and ordinary matrix trace.
Closing under sums and retracts supplies duals for the whole generated
target.
\end{proof}

\begin{theorem}[Temperley--Lieb and Fibonacci are separate quotient branches]
{\normalfont\statusProved}\quad The Temperley--Lieb branch requires its
named tensor ideal, cups, caps, dagger, and Jones--Markov trace.  The faithful
coefficient trace cannot descend through the nonzero quotient ideal.  At
$q=1$ and $d=2$, coefficient-trace weights $(1/2,1/2)$ on the symmetric and
alternating projections become physical weights $(3/4,1/4)$.  The Fibonacci
branch further requires the displayed order-five parameter, Jones trace
radical quotient, additive Karoubi completion, and even subcategory.
\end{theorem}
\begin{scope}
Neither the $d=2$ spin tower nor Fibonacci is identified with the full
positive-real endpoint.  Fibonacci is not obtained from any real $q>0$ in
the displayed loop normalization.
\end{scope}
\provenance{F1-LIM-TL,D1237,D1238,D1239}{theory/sidequests/f1-limit/fusion.md}{theory/checks/f1\_operational\_check.py (H10, finite TL/trace sample; no root-of-unity gate)}{theory/verdicts/f1-limit-adjudication.md}
\begin{proof}
The element $f_i$ is the spectral projection for the $(-1)$-eigenvalue of
$T_i$.  The quotient relation gives
\[
 f_if_{i+1}f_i=\frac{q}{(q+1)^2}f_i=\delta^{-2}f_i.
\]
Yet
\[
 \tau_n(f_i)=\frac{q}{q+1},
 \qquad \operatorname{tr}_J(f_i)=\delta^{-2}
   =\frac{q}{(q+1)^2}.
\]
If the faithful coefficient trace descended through a nonzero quotient
kernel, a nonzero $z$ in that kernel would satisfy both
$\tau_n(z^*z)>0$ and $\tau_n(z^*z)=0$, a contradiction.  At $q=1$, the
two projections $(1\pm(12))/2$ have coefficient weights $1/2,1/2$, whereas
their ranks on $(\mathbb C^2)^{\otimes2}$ are $3$ and $1$.

On the positive-real branch,
$\delta=\sqrt q+q^{-1/2}\geq2$.  The Fibonacci value
$2\cos(\pi/5)=(1+\sqrt5)/2$ is smaller than two, so it requires the complex
root and sign convention in the definition.  At that root, the raw
Temperley--Lieb algebra is semisimplified by its Jones trace radical; the
even surviving labels give the stated Fibonacci fusion rule.  The three
traces---coefficient, finite tensor, and Jones---therefore belong to three
different comparison maps and cannot be silently identified.
\end{proof}
